% ======================================================================
% Example section for the thesis: hardware and software environment.
% Same style as RESULTS_REPORT.tex (prose, no ";", no em dash).
% Drop this in as its own section or fold it into the Methodology.
% ======================================================================

\section{Experimental setup}
\label{sec:setup}

All experiments in this chapter were run on a single desktop workstation.
The purpose of describing it in detail is reproducibility. The two local
vision language models and the YOLO baselines are sensitive to the
inference backend and to the amount of graphics memory available, so the
numbers reported here, in particular the latencies, are tied to this
specific configuration. The GPT-4o results are the exception, because that
model was accessed through a hosted interface and its response times
depend on the provider rather than on the local machine.

\subsection{Hardware}

The workstation is built around an AMD Ryzen 9 9950X processor with
sixteen physical cores and thirty two threads, paired with 128 gigabytes
of DDR5 memory running at 5200 mega transfers per second. Local model
inference used a single NVIDIA GeForce RTX 5070 Ti graphics card with
sixteen gigabytes of GDDR7 memory. The integrated graphics of the Ryzen
processor were present but were not used for any inference. Storage was a
two terabyte WD\_BLACK SN770 NVMe solid state drive, which matters only
because every test reads video frames from disk, and the mainboard was a
Gigabyte X870 GAMING WIFI6. Table~\ref{tab:setup-hw} lists the
configuration in full.

\begin{table}[h]
\centering
\small
\setlength{\tabcolsep}{8pt}
\begin{tabular}{@{}l l@{}}
\toprule
Component & Specification \\
\midrule
CPU              & AMD Ryzen 9 9950X, 16 cores, 32 threads \\
GPU             & NVIDIA GeForce RTX 5070 Ti, 16 GB GDDR7 \\
System memory   & 128 GB DDR5-5200 (4 $\times$ 32 GB) \\
Storage         & WD\_BLACK SN770 NVMe SSD, 2 TB \\
Mainboard       & Gigabyte X870 GAMING WIFI6 (AMD X870) \\
Operating system & Windows 11 Pro, version 24H2, build 26200 \\
GPU driver      & NVIDIA 591.86, CUDA 13.1 \\
\bottomrule
\end{tabular}
\caption{Workstation used for all experiments. The local vision language
models and the YOLO baselines ran on the RTX 5070 Ti. GPT-4o ran remotely
through the OpenAI interface.}
\label{tab:setup-hw}
\end{table}
\FloatBarrier

\subsection{Software and model serving}

The evaluation scripts are written in Python 3.14. Video decoding and
frame extraction use OpenCV 5.0. The YOLO baselines use Ultralytics 8.4
with the stock \texttt{yolo11n} and \texttt{yolo11n-pose} weights, run on
the graphics card. The two local vision language models are served by
Ollama 0.24. The Qwen3-VL model is the eight billion parameter build,
about six gigabytes on disk, and the Llama 3.2 Vision model is the eleven
billion parameter build quantised to four bits, about eight gigabytes on
disk. Both fit comfortably in the sixteen gigabytes of graphics memory,
so neither spilled to system memory during the runs. GPT-4o was called
through the OpenAI Python client version 2.4 against the \texttt{gpt-4o}
model. The same model was used as the Stage~3 judge.

A few backend settings were fixed for the whole battery and are worth
recording because they affect the outputs. The Stage~3 judge was called
with temperature zero so that its grades are stable across repeated runs
on identical text. Qwen3-VL is a reasoning model, and it was forced into a
non-thinking mode with an explicit instruction to answer tersely and with
its output token budget raised to two thousand tokens, because with the
default budget its answer was sometimes cut off before it reached the
verdict line. Llama 3.2 Vision accepts only one image per call, so its
pipeline captions each frame separately and then asks the model to
synthesise a verdict from the captions, which is the reason it makes one
model call per frame and is the slowest of the three. The OpenAI API key
is read from a local environment file and is never sent anywhere other
than to the OpenAI interface.

\subsection{Reproducibility}

The frame sampling is deterministic. Stage~1 selects nine indices spread
evenly across the analysed span, and Stage~2 selects a dense set around
the anchor frame using a fixed stride and a fixed count, so two runs of
the same test on the same video see the same frames. The variation between
repeated runs, which the anomaly tests deliberately expose by repeating
every clip ten times, comes from the models themselves and not from the
sampling. The YOLO baselines are fully deterministic, which is why they
were run once per clip and the row was duplicated to line up with the ten
language model runs rather than being executed ten times. Every run writes
its full Stage~1 and Stage~2 text to disk, and the summary tables in this
chapter are regenerated from four comma separated files by the analysis
scripts, so the path from raw output to reported number can be retraced.

\begin{table}[h]
\centering
\small
\setlength{\tabcolsep}{8pt}
\begin{tabular}{@{}l l@{}}
\toprule
Software & Version and role \\
\midrule
Python              & 3.14, evaluation and analysis scripts \\
OpenCV             & 5.0, video decoding and frame extraction \\
Ultralytics YOLO   & 8.4, \texttt{yolo11n} and \texttt{yolo11n-pose} baselines \\
Ollama             & 0.24, local serving of Qwen3-VL and Llama Vision \\
Qwen3-VL           & 8B build, non-thinking mode, 2000 output tokens \\
Llama 3.2 Vision   & 11B build, 4-bit quantised, one image per call \\
OpenAI client      & 2.4, \texttt{gpt-4o} for detection and for the Stage~3 judge \\
\bottomrule
\end{tabular}
\caption{Software stack. Local models ran on the GPU through Ollama and
Ultralytics. GPT-4o and the Stage~3 judge ran through the OpenAI
interface.}
\label{tab:setup-sw}
\end{table}
\FloatBarrier
